\section{The original generated pair on whole base opens}

\begin{lemma}[The actual base change of the same generated pair]
\label{lem:kn-nadel-nef-big-same-pair-base-open}
\lean{EnvelopeProperty.Vanishing.NadelNefBigGlobalPairBaseOpen.source_isPullback,EnvelopeProperty.Vanishing.NadelNefBigGlobalPairBaseOpen.baseTarget,EnvelopeProperty.Vanishing.NadelNefBigGlobalPairBaseOpen.baseTarget_isPullback,EnvelopeProperty.Vanishing.NadelNefBigGlobalPairBaseOpen.baseTargetProjectiveIso,EnvelopeProperty.Vanishing.NadelNefBigGlobalPairBaseOpen.baseTargetProjectiveIso_hom_coeff,EnvelopeProperty.Vanishing.NadelNefBigGlobalPairBaseOpen.baseTargetProjectiveIso_hom_structure,EnvelopeProperty.Vanishing.NadelNefBigGlobalPairBaseOpen.basePair,EnvelopeProperty.Vanishing.NadelNefBigGlobalPairBaseOpen.basePair_fst,EnvelopeProperty.Vanishing.NadelNefBigGlobalPairBaseOpen.basePair_snd,EnvelopeProperty.Vanishing.NadelNefBigGlobalPairBaseOpen.basePair_isPullback,EnvelopeProperty.Vanishing.NadelNefBigGlobalPairBaseOpen.baseProjectivePair,EnvelopeProperty.Vanishing.NadelNefBigGlobalPairBaseOpen.baseProjectivePair_coeff,EnvelopeProperty.Vanishing.NadelNefBigGlobalPairBaseOpen.baseProjectivePair_structure}
\leanok
\uses{lem:kn-nadel-nef-big-same-global-pair-proper}
Let $D$ be an ideal on a scheme $X$ over a field $k$, let $B$ be an
invertible line, and let $G$ be a finite generating family of $D\otimes B$.
Write $\pi:Y\to X$ for the original blowup and
$F=(\pi,q):Y\to T=X\times_k\mathbf P_k^s$ for the same generated pair.
For an affine open $U$ of $X$, put $A=\Gamma(X,U)$ and
$Y_U=\operatorname{Bl}_{D(U)}\operatorname{Spec}A$.
There is a canonical actual Cartesian square
\[
\begin{array}{ccc}
Y_U&\xrightarrow{F_U}&T\times_X\operatorname{Spec}A\\
\downarrow&&\downarrow\\
Y&\xrightarrow{F}&T.
\end{array}
\]
The upper target is canonically $\mathbf P_A^s$.  Under this
identification the structure map of $F_U$ is the original local blowup
projection, and its composite with $\mathbf P_A^s\to\mathbf P_k^s$
is precisely the original $q$ composed with $Y_U\to Y$.
\end{lemma}
\begin{proof}
The actual affine blowup square is Cartesian over the open subscheme
$U$.  Composing its upper projection with $U\simeq\operatorname{Spec}A$
gives the same Cartesian square over $\operatorname{Spec}A\to X$.
The original first projection $T\to X$ and this base morphism have their
canonical pullback.  The original blowup inclusion followed by $F$ and
the original local projection satisfy its defining equality, since
$F$ has first projection $\pi$.  Their actual pullback lift is $F_U$.
Cancel the right-hand Cartesian square from the original blowup square;
the remaining square is the asserted Cartesian square for $F_U$.

Paste the defining square of $T\times_X\operatorname{Spec}A$ with
the square defining $T$.  Its lower coefficient morphism is exactly
$\operatorname{Spec}A\to X\to\operatorname{Spec}k$.
The native projective coefficient square is Cartesian over this same
coefficient morphism.  Uniqueness of the two pullbacks therefore gives
the canonical isomorphism to $\mathbf P_A^s$.  The two projection
equations of that isomorphism and of the original pullback lift give
the two stated identities.
\end{proof}

\begin{lemma}[Composition of the actual ideal--line generator transport]
\label{lem:kn-nadel-nef-big-comap-tensor-composition}
\notready
\uses{lem:kn-nadel-nef-big-global-frame-untwist}
For actual morphisms $f:V\to U$ and $g:U\to X$, the canonical map
$(gf)^*(D\otimes B)\to D\mathcal O_V\otimes(gf)^*B$ agrees with
the two successive canonical ideal--line maps, through the original
pullback composition isomorphisms and the literal equality of the two
geometric ideal comaps.  In particular this comparison preserves every
original basis map of a generating family.
\end{lemma}
\begin{proof}
Factor each map into the original inverse tensorator followed by the
canonical ideal-module comap tensored with the identity of the line.
Naturality of the inverse tensorator moves the first ideal-module map
through the second tensorator.  The two ideal-module maps then compose
to the original composite ideal-module map, followed by the canonical
isomorphism from equality of the geometric ideal comaps.  The inverse
tensorators compose through the original pullback composition
isomorphism because that isomorphism is monoidal.  These are equalities
of the whole module-sheaf maps.  Apply the equality to each original
basis map, using composition of the original unit pullback maps, to
obtain the last assertion.
\end{proof}

\begin{lemma}[Identification with the same framed affine map]
\label{lem:kn-nadel-nef-big-base-open-same-framed-map}
\notready
\uses{lem:kn-nadel-nef-big-same-pair-base-open,lem:kn-nadel-nef-big-comap-tensor-composition,lem:kn-nadel-nef-big-framed-affine-pair-closed}
If $B$ has its actual frame on $U$, the map $F_U$, under the preceding
canonical target identification, is the original framed-affine map of
the actual restricted source family $G$.  Hence this same $F_U$ is a
closed immersion on the entire inverse image of the base open.
\end{lemma}
\begin{proof}
Restrict $G$ by the actual ideal--line comap map.  The original affine
ideal-comap equality identifies its ideal factor with $\widetilde{D(U)}$;
apply this equality as a module isomorphism to the generating family,
retaining each original index.  Pull the actual frame of $B$ to
$\operatorname{Spec}A$.  The preceding composition equality identifies
the original transported family on $Y$, restricted to $Y_U$, with the
successively transported family from this same affine source family.
The canonical exceptional-ideal equality and the original pullback
composition maps identify their target lines.

On every original Rees pivot chart, composing the basis-map equality
with the actual pivot frame identifies every projective coefficient.
The added zero coefficient and the original coefficient homomorphism
are preserved as well.  Composition of affine projective tuples with
the original projective coefficient map is the tuple over the composite
coefficient homomorphism.  The original source chart cover therefore
gives equality of the two whole projective maps.  The preceding two
projection identities identify this projective map with $F_U$ under
the canonical target isomorphism.  The framed-affine theorem makes it
closed, hence the same is true before the target isomorphism.
\end{proof}

This comparison concerns the same original generated pair.  Target
locality on a framed affine cover will give its global monomorphism;
properness of the already constructed pair then gives its closedness.
No numerical nef or big conclusion follows from these statements.
